\documentclass[10pt,a4paper]{article}

\usepackage[final]{bubtcompetition}

\usepackage{graphicx}
\usepackage{booktabs}
\usepackage{multirow}
\usepackage{amsmath}
\usepackage{amssymb}
\usepackage{url}
\usepackage[hidelinks]{hyperref}


% =========================================================
% PAPER INFORMATION
% =========================================================

\bubtpaperid{BNSC-026}

\bubtteamname{Team Phoenix}

\bubttaskname{Fine-Grained Bengali Toxicity Understanding}

\bubtrunningauthor{Mithun et al.}

\bubtrunningtitle{
A Retrieval-Augmented Transformer for Bengali Toxicity Understanding
}


% =========================================================
% TITLE
% =========================================================

\title{
Team Phoenix at the BUBT Bengali NLP Safety Challenge 2026:
A Retrieval-Augmented Transformer for Fine-Grained Bengali Toxicity Understanding
}


% =========================================================
% AUTHORS
% =========================================================

\author{
Md. Mithun Hossain \quad
Author Two \quad
Author Three
\\
Department of Computer Science and Engineering
\\
Bangladesh University of Business and Technology (BUBT)
\\
Dhaka, Bangladesh
\\
\texttt{mithun@bubt.edu.bd}
}


% =========================================================
% DOCUMENT
% =========================================================

\begin{document}


\maketitle

\thispagestyle{bubtfirstpage}


% =========================================================
% ABSTRACT
% =========================================================

\begin{abstract}

This paper describes our system submitted to the BUBT
Bengali NLP Safety Challenge 2026. The competition focuses
on fine-grained toxicity understanding in Bengali. Our
system combines a pretrained Bengali language model with
task-specific prediction heads and additional representation
learning components. We describe the system architecture,
training procedure, experimental setup, and official
evaluation results. We further provide ablation experiments
and error analysis to better understand the behavior of the
proposed approach.

\end{abstract}


% =========================================================
% INTRODUCTION
% =========================================================

\section{Introduction}

The BUBT Bengali NLP Safety Challenge 2026 focuses on
fine-grained toxicity understanding in Bengali. Participating
teams are required to identify whether a Bengali text is toxic
and determine how the toxicity is expressed.

This system description paper presents the methodology,
experimental configuration, results, and analysis of our
submitted system.


% =========================================================
% TASK DESCRIPTION
% =========================================================

\section{Task Description}

The task contains seven binary labels:

\begin{itemize}

    \item \texttt{is\_toxic}

    \item \texttt{implicit\_toxicity}

    \item \texttt{sarcasm}

    \item \texttt{contains\_profanity}

    \item \texttt{threatening}

    \item \texttt{identity\_attack}

    \item \texttt{sexual\_content}

\end{itemize}


% =========================================================
% DATASET
% =========================================================

\section{Dataset}

Table~\ref{tab:data} summarizes the dataset splits.




% =========================================================
% METHODOLOGY
% =========================================================

\section{Methodology}

Describe the submitted system architecture in sufficient
detail for reproducibility.


\subsection{Text Encoder}

Describe the pretrained language model, tokenizer,
maximum sequence length, and fine-tuning strategy.


\subsection{Prediction Layer}

Describe the multilabel prediction architecture and
activation functions.


% =========================================================
% TRAINING OBJECTIVE
% =========================================================

\subsection{Training Objective}

We optimize the model using binary cross-entropy (BCE)
loss for multi-label classification. Given $C$ labels,
the training objective is defined as

\begin{equation}
\mathcal{L}_{\mathrm{BCE}}
=
-\frac{1}{C}
\sum_{c=1}^{C}
\left[
y_c \log(p_c)
+
(1-y_c)\log(1-p_c)
\right].
\label{eq:bce}
\end{equation}

Here, $y_c \in \{0,1\}$ denotes the ground-truth label
for class $c$, while $p_c \in [0,1]$ denotes the predicted
probability for that class.


% =========================================================
% EXPERIMENTAL SETUP
% =========================================================

\section{Experimental Setup}

Report training hyperparameters, optimizer, learning rate,
batch size, number of epochs, random seeds, hardware, and
software libraries.


% =========================================================
% RESULTS
% =========================================================

\section{Results}

\begin{table}[t]

\centering

\small

\begin{tabular}{lc}

\toprule

System & Macro F1 \\

\midrule

Baseline & 0.XX \\

Submitted System & \textbf{0.XX} \\

\bottomrule

\end{tabular}

\caption{
Official evaluation results.
}

\label{tab:results}

\end{table}


% =========================================================
% ERROR ANALYSIS
% =========================================================

\section{Error Analysis}

Discuss common failure patterns such as implicit toxicity,
sarcasm, code-mixing, spelling variation, and contextual
ambiguity.


% =========================================================
% LIMITATIONS
% =========================================================

\section{Limitations}

Discuss limitations of the proposed system, dataset,
evaluation procedure, and generalization capability.





% =========================================================
% CONCLUSION
% =========================================================

\section{Conclusion}

Summarize the submitted system and the main findings.


% =========================================================
% REFERENCES
% =========================================================

\bibliographystyle{ieeetr}

\bibliography{references}


\end{document}